\section{A leading neutral-scalar interaction}\label{sec:scalar}

\subsection{Partial reduction on a regular patch}
Fix a relatively compact generic patch with $a,b>0$, and allow
the spectator bundle and K3 geometry to vary in small patches.
Write $z=(z_E,z_B)$ for bundle coordinates, $y$ for geometric
hypermultiplet coordinates, and $\chi^I$ for locally unwrapped
periodic two-form coordinates. The tensor multiplet and gravity
remain explicit. At energies below the smallest Higgs mass we
eliminate the three massive vector multiplets as well as the
gapped KK modes. We retain all unbroken gauge and charged matter
fields. The reduced scalar kinetic metric is a smooth gauged
target near this vacuum; it is not the possibly singular vacuum
quotient by the unbroken group.

The nondegeneracy needed for this partial reduction can be checked
in the local model. Let $D:\mathfrak{su}(2)\to T_xV$ be the
broken orbit map. At moment level zero its quaternionic images
obey
\begin{equation}\label{eq:normals}
 \langle I_\alpha D\xi,I_\beta D\zeta\rangle
                 =\delta_{\alpha\beta}Q_x(\xi,\zeta).
\end{equation}
The symplectic pairings of orbit vectors vanish by moment-map
equivariance. Since $D$ is injective, the three moment equations
have an invertible derivative on these nine real normal directions.
Their norm-square potential has positive Hessian there. The three
orbit directions are Goldstone directions; a local gauge slice
removes them. This accounts for the twelve real scalar directions
in the three massive vector multiplets.

At the actual ASD background the same construction has an elliptic
description. If $H$ is its full residual compact gauge group and
$\mathfrak h$ its parallel algebra, the hyperk\"ahler Weitzenb\"ock
formula identifies the ASD cokernel with
$\mathfrak h\otimes\RR^3$. On an $H$-invariant Coulomb slice the
complement of this cokernel has a bounded elliptic right inverse;
its finite kernel is the full retained ASD deformation space.
For an off-shell effective action solve the heavy
Euler--Lagrange equations on this complement. Their leading
linearization is the positive ASD Hessian. The implicit-function
construction gives a smooth $H$-equivariant solution in all retained
variables, with the unbroken moment maps left as a potential.
It agrees with the projected ASD construction on neutral ASD vacua;
these equations need not agree away from that locus. The metric,
two-form and fermionic complementary operators are those specified
in section~\ref{sec:vectors}. All inverse bounds here belong to
the fixed regular patch, and can deteriorate when its Higgs gap
vanishes.

The perturbative expansion below means the finite jet of this
tree action, order by order in $\eta$ and external derivatives.
The coefficient functions and the parameter derivatives actually
used are smooth on the fixed patch. The notation $\Oformal(\eta^j)$
specifies the absence of lower coefficients; it does not bound an
unconstructed infinite remainder at finite $\eta$.

\subsection{Mixed constraints and the vertical kinetic coefficient}
For a pure bundle velocity at fixed leading geometry, use a
parameter gauge connection $\Lambda_i$ and set
\begin{equation}\label{eq:horizontal-ASD}
 a_i=\partial_i\mathcal A-d_{\mathcal A}\Lambda_i,\qquad
 d_{\mathcal A}^*a_i=0,\qquad d_{\mathcal A}^+a_i=0.
\end{equation}
The external-index gauge equation is precisely the first condition;
$\mathcal A_\mu=\Lambda_i\partial_\mu z^i$ gives
$F_{\mu m}=a_{i,m}\partial_\mu z^i$. For stable determinant-fixed
$E,B$, the adjoint Laplacian is invertible off its retained central
kernel. Smooth slice dependence and elliptic differentiation control
$\Lambda_i$ and any fixed finite number of derivatives on the
patch. K3 geometric variations lift to nearby ASD connections:
the trace-free stable ASD cokernel is zero, since its harmonic
self-dual adjoint forms have parallel coefficients. The neutral
central equations are compatible with the fixed trivial determinants.

Put $f_i=\partial_i f$ and
$(j_i)_m=\BF(F_{mn},a_i{}^n)$. For the convention
$G_{\mu m}=\eta v_{i,m}\partial_\mu z^i$ the mixed Einstein
equation to the required order is
\begin{equation}\label{eq:Einsteinconstraint}
 \Delta_1v_i=df_i+j_i,\qquad d^*v_i=0.
\end{equation}
Its coefficients follow directly from
\eqref{eq:heterotic-action}: the conformal internal metric contributes
$-3\eta\,df_i/2$, the dilaton contributes $+\eta\,df_i$,
the mixed metric contributes $\eta d^*dv_i/2$, and the gauge
stress contributes $-\eta j_i/2$. The $H$ stress starts at
$\eta^2$; the mixed $R^2$ source for a pure bundle velocity does
also, because the leading tangent connection is independent of $z$.

The longitudinal equation is essential. Differentiating
\eqref{eq:Poisson}, using the Yang--Mills equation and trace invariance,
gives
\begin{equation}\label{eq:longitudinal}
 \Delta_0f_i=\tfrac12\BF(F_{mn},(d_{\mathcal A}a_i)^{mn})
 =\BF(F_{mn},D^ma_i{}^n)=-d^*j_i.
\end{equation}
Thus $df_i+j_i$ is co-closed. Since $H^1(S,\RR)=0$,
\eqref{eq:Einsteinconstraint} has the controlled solution
$v_i=\Delta_1^{-1}(df_i+j_i)$.

For the two-form constraint, write $H_S=\eta h_1$ and
$H_{\mu mn}=\eta(\mathcal H_i)_{mn}\partial_\mu z^i$.
The mixed Bianchi identity and the external-index two-form equation
give
\begin{equation}\label{eq:Bconstraint}
 d\mathcal H_i=\partial_i h_1+\BF(a_i\wedge F)=:S_i,
 \qquad d^*\mathcal H_i=0.
\end{equation}
The source is closed: the derivative of the internal Bianchi
identity is $d\partial_i h_1=-\BF(d_{\mathcal A}a_i\wedge F)$,
and the second term cancels it using $d_{\mathcal A}F=0$.
Since $H^3(S,\RR)=0$, with $G_3$ the ordinary Green operator,
\begin{equation}\label{eq:Bsolution}
 \mathcal H_i=d^*G_3S_i+h_i,
\end{equation}
where $h_i$ is harmonic. These harmonic components describe the
retained periodic two-form coordinates and their local connection;
they are not inverted. Their parameter-space curvature is the
harmonic projection of the universal Chern--Weil form, including
the parameter legs of the universal gauge connection. The universal
Bianchi identity makes it closed. On a contractible patch a smooth
connection potential is obtained by the Poincar\'e homotopy operator.
This accounts for parameter integrability; setting every $h_i$ to
zero independently would not do so. All computed coefficients
and their finite parameter derivatives are bounded on the patch.

The action now determines the leading bundle tensor, not merely its
expected scale. The mixed Yang--Mills term is order $\eta$ and gives
the $L^2$ metric of \eqref{eq:horizontal-ASD}. The leading metric and
dilaton are independent of $z$; their bundle dependence and metric
compensator start at order $\eta$, so their pure-bundle kinetic
contribution starts at $\eta^2$. The two-form constraint has the
same order, with its periodic-coordinate connection of order one.
The leading tangent connection is independent of $z$; its
variation through the corrected metric and torsion starts at
$\eta$. The curvature-squared and Lorentz Chern--Simons terms
therefore do not add an order-$\eta$ pure-bundle tensor. They are
not omitted from base and higher-order coefficients.

If $a_i=\beta_i-\beta_i^*$ with $\beta_i$ its $(0,1)$ part,
the real one-form norm is twice the Dolbeault norm. Restoring the
normalization \eqref{eq:YMaction} gives
\begin{equation}\label{eq:physicalmetric}
 S_{\rm scalar}=-\frac12\int C_B\gKS_{ij}
               \partial_\mu z_E^i\partial^\mu z_E^j+\cdots,
 \qquad C_B=\frac{2L^2}{g_{10}^2}.
\end{equation}
This is a block of the full kinetic tensor, not yet its intrinsic
or ambient curvature.

\subsection{The other neutral fields and ambient curvature}
After a common six-dimensional normalization, the neutral metric
has the following finite-order structure:
\begin{align}\label{eq:blockmetric}
 G_\eta={}&g_{AB}(y)dy^A dy^B
 +\eta h_{ij}(y,z)(dz^i+C_A^i dy^A)(dz^j+C_B^jdy^B)\nonumber\\
 &+\eta^2 k_{IJ}(y)(d\chi^I+V^I)(d\chi^J+V^J)+E_\eta,
 \\
 h(y,z)={}&V_0^{-1}\bigl(g^{\rm KS}_{E,y}(z_E)
                               \oplus g^{\rm KS}_{B,y}(z_B)\bigr).
 \label{eq:productbundle}
\end{align}
The leading blocks $g,h,k$ are positive and smooth. The connection
$V$ has both $dy$ and $dz$ components. Perturbative two-form shift
symmetry makes the coefficient functions independent of $\chi$.
With two parameter derivatives, the formal remainder orders are
\begin{equation}\label{eq:blockorders}
 \begin{gathered}
 E_{yy}=\Oformal(\eta),\quad
 E_{yz},E_{zz}=\Oformal(\eta^2),\\
 E_{y\chi},E_{z\chi},E_{\chi\chi}=\Oformal(\eta^3).
 \end{gathered}
\end{equation}
The leading $y,z$ cross term is included in $C$.
The leading two-form Hodge metric depends only on $y$; its
first correction is one order higher. These orders follow from
the constraints and action count above and agree with the
geometry--bundle--two-form hierarchy of \cite{alpv}. The Lorentz
Chern--Simons and curvature-squared contributions identified there
are included at their required orders.

We spell out the curvature consequence because the small axion
block cannot be treated as a bounded spectator. In the coframe
of \eqref{eq:blockmetric} the scales are $(1,\eta,\eta^2)$.
In coordinates the inverse orders are
\begin{gather*}
 G^{yy}=O(1),\quad G^{yz}=O(1),\quad
 G^{zz}=\eta^{-1}h^{-1}+O(1),\\
 G^{z\chi}=O(\eta^{-1}),\quad G^{y\chi}=O(1),\quad
 G^{\chi\chi}=O(\eta^{-2}).
\end{gather*}
For lower-index Christoffel symbols with two bundle arguments,
\[
 \Gamma_{z,zz}=\eta\Gamma(h)_{z,zz}+O(\eta^2),\quad
 \Gamma_{y,zz}=O(\eta),\quad \Gamma_{\chi,zz}=O(\eta^2).
\]
The last estimate uses shift independence. In the coordinate
curvature formula the second-derivative terms and the $zz$
Christoffel contraction give $\eta R(h)$; all other contractions
have order $\eta^2$ or higher, including the one with
$G^{\chi\chi}$. Consequently
\begin{equation}\label{eq:ambientleading}
 R(G_\eta)_{EEEE}=\frac\eta{V_0}R(g_E^{\rm KS})
                                  +\Oformal(\eta^2).
\end{equation}
The leading product in \eqref{eq:productbundle} is essential for
the pure-$E$ coefficient. The spectator-bundle moduli have not
been declared constant. A nonsingular tensor-field block can
also be retained; its four-$E$ contributions start one order
later by the same calculation. In product-frame units the leading
sectional curvature is therefore $C_B^{-1}K_{\gKS}$.

It remains to compare neutral curvature with the ambient scalar
target including charged fields. Every additional hypermultiplet
in table~\ref{tab:matter} is charged under the residual compact
group. The neutral fields are exactly the $E,B$ moduli and the
geometric/two-form fields. The equivariant partial reduction
constructed above gives a smooth metric on which the full residual
group acts by isometries. Its fixed component is totally geodesic:
a geodesic with fixed initial point and fixed tangent remains
fixed by uniqueness. Thus its four-neutral curvature is the
ambient scalar curvature. The residual global gauge action,
including shared centers, is retained; no further finite or
continuous quotient is taken. This argument does not say that
the $E$-only slice is totally geodesic among the neutral fields;
their effects have instead been included in
\eqref{eq:blockmetric}--\eqref{eq:ambientleading}.

\subsection{The interaction and its order of limits}
Combining \eqref{eq:physicalmetric}, \eqref{eq:ambientleading}
and Proposition~\ref{prop:curvature} proves
\eqref{eq:intro-curvature}. In physical scalar normal coordinates
with the convention assigning positive curvature to a sphere,
$G_{ij}=\delta_{ij}-\tfrac13R_{ikjl}\phi^k\phi^l+\cdots$.
The two-polarization part of the Lorentzian action
$-\tfrac12G_{ij}\partial\phi^i\partial\phi^j$ is
\begin{equation}\label{eq:interaction}
 \mathcal L_4=\frac{\mathcal K_{\rm phys}^{\rm lead}}6
 \left[(\phi^v)^2(\partial\phi^u)^2
       +(\phi^u)^2(\partial\phi^v)^2
       -2\phi^u\phi^v\partial\phi^u\cdot\partial\phi^v\right]
\end{equation}
at the calculated large-volume order. It is a nonzero
curvature-controlled scalar coefficient in the covariant tree
amplitude dictionary \cite{cheung}. Neutral exchange contributions
that appear in other coordinates assemble into this coefficient;
universal graviton-exchange poles are separate.

The statement first computes the leading tree large-volume
coefficient on a fixed regular patch, then takes
$a=\varepsilon$, $b=\kappa\varepsilon$ with fixed $\kappa>0$.
The geometric relative error is $O_\kappa(\varepsilon^2)$.
The derivative expansion requires
$E_{\rm probe}\ll\kappa\varepsilon/L$, with its next parameter
$E_{\rm probe}^2/m_{\min}^2$, and perturbative interactions in
canonical normalization. Taking the probe energy sufficiently
small and the ten-dimensional coupling weak leaves a nonempty
regime. Keeping a six-dimensional dilaton fixed as volume grows
need not keep that ten-dimensional coupling small, so these
limits are specified separately. We have not bounded higher
physical coefficients as $\varepsilon\to0$ or proved a uniform
simultaneous finite-$\eta$ curvature remainder.
